\begin{lemma}
For every partial coloring $\varphi$ of the original edge type and every
$x\in V^{\circ}(G)$, the colors present at $x$ in $G^{\circ}$ are exactly
the colors present at its underlying vertex in $G$.  Consequently the two
missing-color sets are equal as well.
\end{lemma}
\begin{proof}
By \noderef{PresentColors}, a color $c$ is present at $x$ precisely when
some edge identity $e$ incident with $x$ receives $c$.  The edge type and
the coloring function are identical in the two graphs, while
\noderef{IncidentCoreIncidence} identifies the two incidence conditions.
Thus the same edge identity witnesses presence of $c$ on either side, which
proves equality of the present-color sets.

By \noderef{MissingColors}, the missing colors are the complement of the
present colors in the same finite palette.  Taking complements of the
equality just proved gives equality of the missing-color sets.
\end{proof}
